\documentclass[uplatex,dvipdfmx]{jlreq}
\usepackage[fleqn,tbtags]{mathtools}
\usepackage{jlreq-deluxe}
\usepackage[noalphabet]{pxchfon} % must be after otf package
\usepackage{stix2} %欧文＆数式フォント
\usepackage[fleqn,tbtags]{mathtools} % 数式関連 (w/ amsmath)
% \usepackage{hira-stix} % ヒラギノフォント＆STIX2 フォント代替定義（Warning回避）
\usepackage{tabularx}

% FBS書く設定
\usepackage{tikz}
\usetikzlibrary{trees,shadows}
\makeatletter
\@ifundefined{kanjiskip}{}{\usepackage{pxpgfmark}}
\makeatother


\begin{document}

\title{第三回 事前課題: 担当箇所の基本設計・詳細設計のたたき台}
\author{25G1059 佐藤涼菜}
\date{2026年4月22日}
\maketitle




\section{基本設計}

\subsection{指揮デバイス}
・本体(図)→加速度センサ？\par
・指揮デバイス振ったらドラムArduinoに接続\par
・ボタン押したら音量，強さ変化\par
(・ボタン押したらBPM変化？おそらくドラムArduinoになりそ)\par





FBSについて示すにょ

\begin{center}
    \begin{tikzpicture}[
        node distance=1.5cm,
        every node/.style={
            draw=black!70!black,
            thick,
            rounded corners=3pt,
            fill=white,
            align=center,
            font=\small\sffamily,
            inner sep=5pt,
            % drop shadow % 影
        },
        edge from parent/.style={
            draw=black!70!black,
            thick
        },
        level 1/.style={sibling distance=4.5cm},
        level 2/.style={sibling distance=2.2cm}
    ]
    
    % ルート
    \node {指揮デバイス}
        child { node {予定管理 \\ \scriptsize 新規作成/編集/削除}
            child { node {日時の設定} }
            child { node {場所の設定} }
        }
        child { node {通知 \\ \scriptsize メール通知}
            child { node {登録確認 \\ メール} }
            child { node {リマインド \\ メール} }
        }
        child { node {予定の共有 \\ \scriptsize アカウント}
            child { node {アカウント \\ 登録} }
            child { node {共有ユーザ \\ の選択} }
        };
    \end{tikzpicture}
    \end{center}




\section{詳細設計}

\subsection{サーボモータ}
    \subsubsection{メトロサーボ}
    BPMに合わせてぶん回す

    \subsubsection{ドラムサーボ}

    \subsection{LED}
    BPMに合わせて光る

    \subsection{LEDマトリックス}
    楽曲の音程に合わせて表示(グラデーション？)

















\begin{thebibliography}{999}

\bibitem{ref:kitai}星野悦子．音楽心理学の研究動向ー音楽で喚起される感情の研究を中心にー．音楽教育学第51-2.(2022)．
\bibitem{ref:eko}Helmut Haas．Über den Einfluss eines Einfachechos auf die Hörsamkeit von Sprache．Acustica, Vol. 1, pp. 49-58．
\bibitem{ref:yamaha}YAMAHA.SYNCROOM公式HP．
\small\url{https://syncroom.yamaha.com/}

\bibitem{ref:adsr}島村．アナログシンセ超入門～その2：EG（エンベロープジェネレーター、ADSR）編.(2026年4月21日閲覧)
\small\url{https://info.shimamura.co.jp/digital/guide/2018/02/122103}

\end{thebibliography}

\end{document}




